\section{Part I: Why Research Is a Different Enterprise}

\subsection{1. The difference between delivering expertise and producing new knowledge}

A consulting team deploys knowledge it already possesses to solve a client's problem. A product team engineers a repeatable solution under constraints. A research team investigates whether a proposition is true, whether a mechanism exists, or whether a method can do something previously unknown or poorly understood. The activities overlap, but their success criteria differ. A technically difficult project is not necessarily research; neither is every new customer feature a scientific contribution. Novelty to the team differs from novelty to the field.

The founder's 1st act is to define the \textbf{research object}: an unknown that, if resolved, will improve understanding or materially change a subsequent decision. ``Build an AI lab'' is not a research object. ``Determine whether a specified evaluation protocol detects a class of failures missed by available baselines, with reproducible uncertainty estimates'' may be. A research enterprise begins with an epistemic commitment, not a building or a job title.

\subsection{2. 4 outputs, 4 standards of success}

Research can produce \textbf{knowledge} (reliable claims or negative results), \textbf{capabilities} (trained people and methods), \textbf{assets} (software, datasets with rights, instruments, patents, validated protocols), and \textbf{options} (new products, partnerships, clinical studies, ventures, future experiments). These outputs may not appear simultaneously. A failed hypothesis can be scientifically valuable while commercially disappointing. A patent may have no market value. A research prototype may create an excellent engineer but no invention.

A founder should name the primary intended output of each project before funding it. Assigning all projects a short-term revenue target rewards low-risk consultancy. Assigning none a financial constraint invites an unbounded subsidy. Good governance acknowledges different scoreboards without abandoning comparative allocation.

\subsection{3. Basic, applied, translational, and commissioned work}

\textbf{Basic research} advances understanding without requiring near-term application. \textbf{Applied research} addresses a practical question using new investigation. \textbf{Translational research} moves a finding toward usable conditions-repeatability, engineering, validation, regulatory evidence, manufacturability, or deployment. \textbf{Commissioned development} often builds a client-specific artifact from known methods; it may contain research but is not inherently research.

These form a landscape, not a prestige hierarchy. In a young organization, a narrow applied question with a clear baseline can be more serious and useful than a grand speculative agenda. In a mature university, protected fundamental inquiry may be essential even when its monetization is unknowable. Fund each according to its uncertainty, timescale, and mission.

\subsection{4. Research is a real option, not a promised product}

An experiment costs money today to change what management knows tomorrow. Let an experiment cost $E$; suppose it may reveal whether a larger $K$ investment is worthwhile. Its value can come from \textbf{avoiding an otherwise bad commitment} as much as from revealing a successful opportunity. The value of information is approximately the expected best decision \textit{with} the information minus the expected best decision \textit{without} it, less the cost of obtaining information. Calculating it requires assumptions about possible states and actions; a spreadsheet does not turn subjective probabilities into facts.

An enterprise often should spend ₹3 lakh to determine that a ₹50-lakh product is unworkable. Conversely, a founder can waste ₹3 lakh on a beautiful experiment that does not alter any future decision. The test is whether plausible results lead to genuinely different actions.

\subsection{5. The asymmetry between failure and misconduct}

A null result, inconclusive result, rejected hypothesis, or failed prototype can be evidence of good research. Manipulated measurements, invented results, hidden conflicts, or misleading presentation are not. The US Office of Research Integrity's formal definition focuses on fabrication, falsification, and plagiarism, while noting that honest error and differences of opinion are not misconduct [R1]. The exact local legal standard depends on jurisdiction and funding agreement.

Research leaders should reward well-designed, candidly reported tests and punish dishonest inference. A lab that penalizes every negative outcome will soon report too many positive ones. A commercial team that treats all results as marketing assets has created pressure precisely where independent truth-seeking is needed.

\subsection{6. Research economics are dominated by the scarce complement}

Cash buys equipment and salaries, but cannot automatically buy an important question, scientific taste, access to valid samples, replication discipline, or acceptance by a serious field. The bottleneck may be a senior experimentalist, permission to use data, patient follow-up, a clinical site, an instrument technician, or protected time. A cheap computer-science study can be infeasible without a legitimate dataset; an expensive wet lab can remain scientifically unproductive without a capable principal investigator.

Thus the economically correct order is \textbf{question → protocol → people and permissions → infrastructure → funding instrument}, revisiting each as evidence appears. Buying infrastructure 1st can reverse causality: the purchased apparatus begins choosing the problems.

\subsection{7. The 3 clocks and their incompatible incentives}

Commercial services operate on customer and cash cycles; academic research on discovery and verification cycles; venture creation on customer adoption and financing cycles. A 12-week enterprise contract may be an excellent way to pay bills but a terrible unit for measuring fundamental research. A scientific group that cannot publish for 3 years may nevertheless be doing responsible work, but must explain why continued funding is rational.

Design separate milestones: \textbf{scientific} (methods, evidence, replication), \textbf{engineering} (performance in realistic conditions), \textbf{commercial} (buyer adoption and contribution), and \textbf{financial} (cash obligations, runway, restrictions). A milestone in 1 dimension does not certify the others.

\subsection{8. The founder's 1st question: what is our distinctive right to investigate?}

An organization should not enter an area solely because it is fashionable. Ask whether it has: unusual questions grounded in real problems; genuinely qualified investigators; access to useful environments; reliable methods; prospective collaborators; the means to supervise and audit work; and enough runway to survive a null result. A new institute's unusual opportunity may be the intersection of practice and pedagogy: observed failures in real AI deployments and in the development of human AI-system judgment.

``Because we teach AI, we should research all of AI'' is not a thesis. ``We repeatedly see a costly failure in applied AI evaluations, have permission to study it, and can test a specific improvement against strong baselines'' is closer to 1.
